% ============================================================
%  Chapter 3: R3 — Local Gauge Invariance
% ============================================================
\chapter{R3 — 局域规范不变性}
\label{ch:r3}

\section{总说：对称性决定力}

R3 是对现代物理学影响最深远的原理之一。它说：\textbf{物理定律在局域的、位置依赖的对称性变换下保持形式不变}。为维持这种不变性，必须引入补偿场——规范场——而这些规范场就是基本相互作用的传递者。

这意味着：\textbf{力的存在不是随意附加的，而是对称性的必然要求}。电磁力、弱力、强力，都源自同一条原理——只是规范群不同。

\begin{principlebox}[R3 — 局域规范不变性]
  物质场的拉格朗日量在任意局域（位置依赖的）规范变换下保持不变。
  为维持这种不变性，必须引入规范场 $A_\mu^a$。
  规范场与物质场的耦合自然产生基本相互作用力。
  规范群的结构完全决定了相互作用的形式。
\end{principlebox}

\textbf{注意}：R3 不指定规范群的具体选择（$U(1)$, $SU(2)$, $SU(3)$ 等）——那是 contingent 事实。
R3 只规定：\textbf{若}存在规范对称性，则必然伴随着规范场和规范相互作用。

% ── 变量定义表 ──
\section{变量定义}

\begin{table}[h]
\centering
\small
\begin{tabular}{p{2.2cm}p{1.8cm}p{1.8cm}p{1.5cm}p{3cm}p{4cm}}
\toprule
\textbf{变量} & \textbf{符号} & \textbf{类型} & \textbf{量纲} & \textbf{定义域} & \textbf{物理含义} \\
\midrule
物质场 & $\psi(x)$ & 依赖表示 & 依赖场类型 & $x \in \mathcal{M}$ & 承载规范荷的物质场 \\
规范参数 & $\alpha^a(x)$ & $\mathbb{R}$ & [1] & 任意光滑函数 & 局域规范变换的角度 \\
规范势 & $A_\mu^a(x)$ & $\mathbb{R}$ & 依赖群 & $\mu \in \{0,1,2,3\}$ & 补偿场；传递力的玻色子 \\
场强张量 & $F_{\mu\nu}^a$ & 反对称 & 依赖群 & $F_{\mu\nu}^a = -F_{\nu\mu}^a$ & 规范场的物理自由度 \\
协变导数 & $D_\mu$ & 算符 & 依赖对象 & $\mu \in \{0,1,2,3\}$ & 规范协变的导数 \\
规范耦合常数 & $g$ & $\mathbb{R}^+$ & [1] & $g > 0$ & 物质场-规范场耦合强度 \\
生成元 & $T^a$ & 矩阵 & [1] & $a = 1,\dots,\dim(G)$ & Lie 代数的基础表示 \\
结构常数 & $f^{abc}$ & $\mathbb{R}$ & [1] & $[T^a,T^b] = if^{abc}T^c$ & 非 Abel 特性度量 \\
\bottomrule
\end{tabular}
\caption{R3 变量定义}
\end{table}

% ── 数学表达式 ──
\section{数学表达式}

\subsection{局域规范变换（U(1) 为例）}
\[
\psi(x) \to e^{i\alpha(x)}\psi(x)
\]

\subsection{协变导数（最小耦合）}
\[
D_\mu \equiv \partial_\mu - i g A_\mu(x)
\]
要求 $D_\mu\psi \to e^{i\alpha} D_\mu\psi$，由此确定规范场变换律：
\[
A_\mu(x) \to A_\mu(x) + \frac{1}{g}\,\partial_\mu\alpha(x)
\]

\subsection{场强张量}
\[
F_{\mu\nu} \equiv \partial_\mu A_\nu - \partial_\nu A_\mu
\]
非 Abel 推广：$F_{\mu\nu}^a \equiv \partial_\mu A_\nu^a - \partial_\nu A_\mu^a + g f^{abc} A_\mu^b A_\nu^c$。

\subsection{规范不变拉格朗日量}
纯规范场（SI 单位，U(1) EM 为例）：
\[
\mathcal{L}_{\text{gauge}} = -\frac{1}{4\mu_0} F_{\mu\nu} F^{\mu\nu}
\]
系数 $1/\mu_0$ 由实验确定——规范对称性本身不固定耦合强度。

耦合到 Dirac 物质场（自然单位 $\hbar=c=1$）：
\[
\mathcal{L} = \bar{\psi}(i\gamma^\mu D_\mu - m)\psi - \frac{1}{4\mu_0}F_{\mu\nu}F^{\mu\nu}
\]

% ── 规范场动力学的确定 ──
\section{规范场动力学的确定}

R3 本身只强制规范场 $A_\mu^a$ 的\textbf{存在性}和\textbf{最小耦合}（$p \to p - gA$）。
动能项 $\mathcal{L}_{\text{gauge}}$ 需要额外原则确定：

\begin{itemize}
  \item Lorentz 不变性 (R1) + 至多二阶导数 $\to$ 唯一形式 $\propto F_{\mu\nu}^a F^{a\mu\nu}$
  \item 比例系数（如 $1/4\mu_0$）由实验确定
  \item 高阶导数项在低能极限下被压低
\end{itemize}

% ── 成立条件与失效边界 ──
\section{成立条件与失效边界}

\begin{limitbox}[R3 的成立条件与失效边界]
  \textbf{成立条件}：
  \begin{itemize}
    \item 可微规范变换：$\alpha(x)$ 充分光滑
    \item 规范群已指定（contingent 事实——$U(1)$, $SU(2)$, $SU(3)$ 等由实验确定）
  \end{itemize}

  \textbf{失效边界}：
  \begin{itemize}
    \item 反常 (Anomaly)：经典规范对称性在量子层面可被破坏
    \item 强耦合区：微扰展开失效，需非微扰方法
    \item Planck 尺度：规范对称性可能 emergent 而非 fundamental
  \end{itemize}
  \textit{注：Higgs 机制不破坏规范对称性——真空期望值破缺的是全局对称性部分，规范对称性本身保持完好。}
\end{limitbox}

% ── 必要性 ──
\section{必要性 (Necessity)}

\begin{table}[h]
\centering
\small
\begin{tabular}{p{5cm}p{7cm}}
\toprule
\textbf{移除/弱化} & \textbf{后果} \\
\midrule
$\alpha(x)$ 局域 $\to$ $\alpha$ 全局常数 & 无规范场 $\to$ 无光子/电磁力；QFT 不可重整 \\
规范不变性 $\to$ 显式破缺 & 非物理自由度不消除；概率不守恒 \\
规范群为 Abel 群（无非 Abel） & 无规范场自耦合 $\to$ 无渐近自由 $\to$ 无 QCD 禁闭 \\
\bottomrule
\end{tabular}
\caption{R3 必要性}
\end{table}

% ── 充分性 ──
\section{充分性 (Sufficiency)}

\textbf{R3 + R1 + R2} 提供：
\begin{itemize}
  \item $U(1)$ 规范群 + 最小作用量 $\to$ Maxwell 方程组
  \item 非 Abel 规范群 $\to$ Yang-Mills 理论
  \item 规范相互作用的普适形式（最小耦合）
  \item 自发对称破缺 + R3 $\to$ Higgs 机制 $\to$ 弱力玻色子质量
\end{itemize}

\section{直觉理解：为什么需要规范场？}

考虑一个自由电子的波函数 $\psi(x)$。如果我们要求物理可观测量在局域相位变换 $\psi \to e^{i\alpha(x)}\psi$ 下不变，导数的普通定义就不够了：

\[
\partial_\mu(e^{i\alpha}\psi) = e^{i\alpha}(\partial_\mu\psi + i\partial_\mu\alpha \cdot \psi) \neq e^{i\alpha}\partial_\mu\psi
\]

多出来的项 $i\partial_\mu\alpha$ 打破了不变性。我们需要一个"修正"的导数 $D_\mu$，使得变换后这多余项被消去。解决办法是引入一个场 $A_\mu$——\textbf{规范场}——让它沿着 $\psi$ 同时变换，恰好抵消多余项：

\[
D_\mu\psi = (\partial_\mu - ieA_\mu)\psi, \quad A_\mu \to A_\mu + \frac{1}{e}\partial_\mu\alpha
\]

\textbf{这就是电磁力的起源}。要求局域相位不变性 $\Rightarrow$ 必须存在光子场 $A_\mu$ $\Rightarrow$ 光子与电子耦合 $\Rightarrow$ 电磁力。

这条原理后来被推广到了更复杂的对称性：$SU(2)$ 产生了弱力（$W^\pm, Z^0$），$SU(3)$ 产生了强力（胶子）。
